\documentclass[aspectratio=169]{beamer}

\usetheme{Madrid}
\usecolortheme{default}
\setbeamertemplate{navigation symbols}{}

\usepackage[T1]{fontenc}
\usepackage[utf8]{inputenc}
\usepackage{booktabs}
\usepackage{array}
\usepackage{tikz}

\title{Economic Supervisor Benchmark}
\subtitle{Dynamic routing vs. a fixed frontier model}
\author{ClawAgents Benchmark}
\date{March 13, 2026}

\begin{document}

\begin{frame}
  \titlepage
\end{frame}

\begin{frame}{Idea}
  \begin{itemize}
    \item Goal: reduce spend by routing each turn to an appropriately capable model instead of always paying for a frontier model.
    \item Tiers in this benchmark:
    \begin{itemize}
      \item Simple: \texttt{gpt-5-nano}
      \item Coding: \texttt{gpt-5-mini}
      \item Reasoning: \texttt{gpt-5.4}
    \end{itemize}
    \item Important: the supervisor does not always try the smallest model first.
    \item It classifies the recent context once, then sends the request directly to one tier.
  \end{itemize}
\end{frame}

\begin{frame}{Routing Logic}
  \begin{block}{Decision rule}
    The supervisor inspects recent messages and chooses one route:
  \end{block}

  \begin{itemize}
    \item \textbf{Reasoning} for distress or deep-reasoning signals
    \item \textbf{Coding} for debugging / code-editing signals
    \item \textbf{Simple} for low-complexity requests
  \end{itemize}

  \vspace{0.5em}

  \begin{alertblock}{Implication}
    This is a \textbf{single-shot router}, not a cheap-first retry loop.
  \end{alertblock}
\end{frame}

\begin{frame}{Benchmark Setup}
  \begin{itemize}
    \item Baseline: fixed \texttt{gpt-5.4}
    \item Supervisor: \texttt{gpt-5-nano} / \texttt{gpt-5-mini} / \texttt{gpt-5.4}
    \item Repeats: 3
    \item Settings: streaming enabled, rethink enabled, trajectory disabled
    \item Workload: 9 repository tasks covering file reads, summaries, shell execution, testing, and design-level analysis
  \end{itemize}

  \vspace{0.5em}

  \begin{table}
    \centering
    \begin{tabular}{lcc}
      \toprule
      Scenario & Total spend & Total tokens \\
      \midrule
      Fixed model & \$3.197865 & 213{,}191 \\
      Economic supervisor & \$2.123240 & 215{,}343 \\
      \bottomrule
    \end{tabular}
  \end{table}
\end{frame}

\begin{frame}{Headline Result}
  \begin{columns}[T]
    \column{0.48\textwidth}
    \begin{block}{Spend}
      \centering
      {\LARGE -33.6\%}

      \vspace{0.3em}
      \$3.197865 $\rightarrow$ \$2.123240
    \end{block}

    \column{0.48\textwidth}
    \begin{block}{Tokens}
      \centering
      {\LARGE +1.01\%}

      \vspace{0.3em}
      213{,}191 $\rightarrow$ 215{,}343
    \end{block}
  \end{columns}

  \vspace{0.8em}

  \begin{alertblock}{Main takeaway}
    The supervisor preserved roughly the same token volume while cutting total spend by about one third.
  \end{alertblock}
\end{frame}

\begin{frame}{Where the Calls Went}
  \begin{columns}[T]
    \column{0.52\textwidth}
    \begin{table}
      \centering
      \begin{tabular}{lrr}
        \toprule
        Tier & Calls & Share \\
        \midrule
        Simple & 12 & 16.4\% \\
        Coding & 27 & 37.0\% \\
        Reasoning & 34 & 46.6\% \\
        \bottomrule
      \end{tabular}
    \end{table}

    \column{0.48\textwidth}
    \begin{table}
      \centering
      \begin{tabular}{lr}
        \toprule
        Model & Spend \\
        \midrule
        \texttt{gpt-5-nano} & \$0.008578 \\
        \texttt{gpt-5-mini} & \$0.122122 \\
        \texttt{gpt-5.4} & \$1.992540 \\
        \bottomrule
      \end{tabular}
    \end{table}
  \end{columns}

  \vspace{0.5em}

  \begin{itemize}
    \item Cheap tiers handled 39 of 73 calls (53.4\%).
    \item Those 39 calls cost only \$0.130700 total.
    \item Most spend still came from reasoning turns on \texttt{gpt-5.4}.
  \end{itemize}
\end{frame}

\begin{frame}{Task-Level Readout}
  \begin{table}
    \centering
    \small
    \begin{tabular}{>{\raggedright\arraybackslash}p{0.49\textwidth}rr}
      \toprule
      Task & Spend delta & Latency delta \\
      \midrule
      Read \texttt{pyproject.toml} version & -90.24\% & +145.43\% \\
      Run \texttt{pwd} & -97.06\% & +141.50\% \\
      Summarize test behavior & -59.37\% & -54.43\% \\
      Run \texttt{pytest} on the supervisor test & -76.16\% & +297.61\% \\
      Explain latency + cost metrics & -55.63\% & -5.85\% \\
      Compare fixed vs. supervisor & +22.26\% & +19.44\% \\
      \bottomrule
    \end{tabular}
  \end{table}

  \vspace{0.4em}

  \begin{itemize}
    \item 8 of 9 benchmark tasks were cheaper under the supervisor.
    \item Only 3 of 9 tasks were faster.
    \item Savings were strongest on simple tool or lookup tasks; latency often increased.
  \end{itemize}
\end{frame}

\begin{frame}{Interpretation}
  \begin{block}{What worked}
    The router successfully offloaded a meaningful share of turns to cheaper models, producing large cost reductions on low-complexity tasks.
  \end{block}

  \begin{block}{What did not improve}
    Lower spend did not reliably reduce latency. More routing overhead, extra iterations, or weaker cheap-tier performance can erase speed gains.
  \end{block}

  \begin{alertblock}{Design lesson}
    Cost and latency must be measured separately. A routing policy can be economically better while still being operationally slower.
  \end{alertblock}
\end{frame}

\begin{frame}{Conclusion}
  \begin{itemize}
    \item The economic supervisor demonstrated the intended idea: spend frontier-model dollars only when the context appears to justify it.
    \item On this benchmark, that reduced spend by 33.6\% with almost no change in total token usage.
    \item The tradeoff was latency: many cheap-route wins were slower overall.
    \item Next step: improve routing quality and measure answer quality, not just cost and speed.
  \end{itemize}
\end{frame}

\end{document}
